% Shared-oracle primitive derivations for joint-continuation.tex, v0.3.
% Place after primitive-inputs.tex. Uses its Phi_256 and E_NMAC conventions.
\section{Shared-oracle primitives for the joint continuation}
\label{app:joint-primitives}

This appendix supplies the primitive reductions used in
Section~\ref{sec:joint-continuation}. They use the same idealized SHA-256
compression oracle throughout. Each experiment begins with independent
oracle, hidden uniform 256-bit input and auxiliary state. All
oracle-dependent preparation is executed within the counted algorithm.
The designated hidden value has only the stated interfaces, and neither
it nor the private random tables are returned to the terminal test.
Ordinary public queries may be quantum; private function queries are
classical and exact repeats are cached.

We use $\Phi_{256}$ and $E_{\rm NMAC}$ from
Appendix~\ref{app:primitive-inputs}, with
\[
 \ell(L)=1+\left\lceil\frac{L+9}{64}\right\rceil,
 \qquad L\le 2^{61}-65.
\]
The HMAC padding includes the consumed 64-byte key pad: its encoded
length is $8(64+|m|)$. Equal padded strings have equal encoded lengths
and hence equal message prefixes within this domain. Thus the required
padding is injective, and $\ell(L)$ bounds the ordinary row compression
work after its two key-pad evaluations.

\subsection{A shared HMAC row and column}
\label{app:joint-row-column}

For a 32-byte value $x$, define the 64-byte blocks
\[
 \begin{split}
 I(x)&=(x\Vert0^{32})\oplus\mathtt{36}^{64},\\
 O(x)&=(x\Vert0^{32})\oplus\mathtt{5c}^{64},\\
 B(x)&=x\Vert P,\qquad
 P=\mathtt{80}\Vert0^{23}\Vert\operatorname{BE64}(768).
 \end{split}
\]
The final 32 bytes of $I(x),O(x),B(x)$ are respectively
$\mathtt{36}^{32},\mathtt{5c}^{32},P$. These three ranges are disjoint,
and each map is injective. Write
\[
 \alpha(a)=h(I(a),\mathrm{IV}),\qquad
 \beta(a)=h(O(a),\mathrm{IV}),\qquad
 \Gamma(a)=(\alpha(a),\beta(a)).
\]
The actual 32-byte-message column is
\begin{equation}
 F(a,R)=h\bigl(B(h(B(R),\alpha(a))),\beta(a)\bigr).
 \label{eq:appendix-column-compression}
\end{equation}
The actual row is the NMAC-form expression
\[
 F(R,m)=\operatorname{Iter}_h\!\left(
 \beta(R),\operatorname{pad}_{64}
 \bigl(\operatorname{Iter}_h(\alpha(R),\operatorname{pad}_{64}(m))\bigr)
 \right),
\]
where $\operatorname{pad}_{64}$ retains the consumed-key-pad length offset.

For completeness, $E_{F_1}(C,q)$ denotes a valid distinguishing bound for
\[
 F_1(u,v)=\widetilde h(v,f(u))
\]
versus a private random function of the complete pair $(u,v)$.
Here $f$ is an independent private random function from 256 to 256 bits;
$\widetilde h$ is a public random function on 512-bit $v$ and 256-bit final
input; there are $C$ classical $F_1$ calls and $q$ ordinary quantum queries.
Hosoyamada and Iwata's Proposition~4~\cite{hmacqrom2021} gives
\[
 E_{F_1}(C,q)=
 O\!\left(\sqrt{\frac{q_*^3}{2^{256}}}\right),
 \qquad q_*=\max\{C,q\}=o(2^{256/3}),
\]
with the published constants and regime retained.

\begin{lemma}[Independent-table shared-root bound]
\label{lem:appendix-joint-hmac}
An admissible operational bound is
\begin{equation}
\begin{split}
 \varepsilon_{\rm Hjoint}(C,T,q,L)\le\min\{1,\;&
 \Phi_{256}(C+2,q+3C+\ell(L)T)\\
 &+E_{\rm NMAC}(T,q+3C,\ell(L))\\
 &+E_{F_1}(C,q+2C)\}.
\end{split}
 \label{eq:appendix-joint-hmac}
\end{equation}
The ideal row and column are independent private functions
$U(m)$ and $V(\Gamma(a))$, independent of the public oracle.
\end{lemma}
\begin{proof}
First isolate every compression call with a raw occurrence of $R$:
\[
 \mathcal F_R(t,s)=h(\tau_t(R),s),
 \qquad \tau_{\rm in}=I,\quad\tau_{\rm out}=O,\quad\tau_{\rm msg}=B.
\]
For a fixed $(t,s)$, each encoded point has probability at most $2^{-256}$.
Different effective pairs encode different points for every $R$:
different tags have disjoint block ranges, and equal tags retain the
chaining state. Cache by the complete pair $(t,s)$, not just by the salt.
This matters when two salts have the same $\alpha$ value.

There are at most $C+2$ distinct private calls: the two root key-pad
evaluations and at most $C$ message-block evaluations at $\alpha(a)$.
The remaining ordinary work is at most $q+3C+\ell(L)T$, including each
column's two public pad states and outer call, and each row's continuation.
Don, Fehr and Huang's Theorem~2~\cite{dfh2022} therefore replaces
$\mathcal F_R$ by a private random function at the first cost in
\eqref{eq:appendix-joint-hmac}.

Its restrictions
\[
 A=\mathcal F({\rm in},\mathrm{IV}),\quad
 D=\mathcal F({\rm out},\mathrm{IV}),\quad
 f(s)=\mathcal F({\rm msg},s)
\]
are independent of one another and $h$. The row is now NMAC with independent
inner and outer keys $A,D$. The simultaneous column is
$h(B(f(\alpha(a))),\beta(a))$. To replace the row, the NMAC reduction
samples its own $f$ table and answers a column with three ordinary
compression calls. Proposition~6 of Hosoyamada and Iwata
\cite{hmacqrom2021}, with the prescribed offset padding, gives the second
term and a private row function $U$ independent of $f,h$.

For the column, reindex the \emph{same} oracle by
\[
 \widetilde h((t,s),y)=h(y\Vert t,s),
 \qquad |t|=|s|=|y|=256.
\]
This is a fixed bijection on all 768 input bits. It is implemented
coherently by wire permutation and one oracle call, not by introducing an
independent hash. The remaining column is exactly
\[
 \widetilde h\bigl((P,\beta(a)),f(\alpha(a))\bigr)
   =F_1\bigl(\alpha(a),(P,\beta(a))\bigr).
\]
Computing its public states costs two ordinary calls per column. The
independent row table $U$ needs no compression queries.
The $F_1$ bound therefore applies with $C$ private calls and at most
$q+2C$ ordinary calls. Its ideal reply is indexed by
$(\alpha(a),P,\beta(a))$, equivalently by $\Gamma(a)$ since $P$ is fixed.
Equal $\Gamma$ values must share a reply; their equality is never assumed
to imply equal salts. Summing the three replacements proves the result.
\end{proof}

The lemma compares against the independent-table ideal. In particular it
does not assert pointwise real equality at an exceptional cross-interface
pair where both classical arguments equal the hidden root. Such guesses
are admitted by the real experiment and covered by the three
distinguishing steps. A different ideal that explicitly repairs this
cross-pair needs its own comparison; it is not used here.

\subsection{The multiple-message row and a single Extract column}

For $\varepsilon_{\rm Hrow}(T,q,L)$, the adversary has an adaptive classical
row oracle. It can choose later messages using earlier replies and its
quantum computation with the common public oracle. Equal messages retrieve
equal values. There is no requirement that the full transcript be fixed
before the first reply, and there is no raw-key disclosure channel.

\begin{lemma}[Row and single-column instances]
\label{lem:appendix-hmac-instances}
The following bounds instantiate the operational games:
\begin{align}
 \varepsilon_{\rm Hrow}(T,q,L)
 &\le\min\{1,\,
 \Phi_{256}(2,q+\ell(L)T)+E_{\rm NMAC}(T,q,\ell(L))\},
 \label{eq:appendix-multirow}\\
 \varepsilon_{\rm Hcol}(q)
 &\le\min\{1,\,
 \Phi_{256}(3,q+3)+E_{F_1}(1,q+2)\}.
 \label{eq:appendix-single-column}
\end{align}
\end{lemma}
\begin{proof}
For the row, only $I(K)$ and $O(K)$ contain the hidden key.
The same disjoint-encoding replacement uses two private calls and at most
$q+\ell(L)T$ ordinary calls. The resulting independent-key NMAC receives
all $T$ adaptive messages, and its replacement is precisely the private
row function. No conversion to a one-transcript game is made.

For one distinct zero-salt column, omit the NMAC step from
Lemma~\ref{lem:appendix-joint-hmac}, put $C=1$, and use its conservative
$C+2=3$ private-call budget. The only visited $\Gamma(0)$ entry gives one
uniform reply independent of $h$ and of the executed prefix. The extra
two private-call slots need not actually be used; retaining this upper
bound agrees with the existing shared-root accounting.
\end{proof}

\subsection{Root initialization at the actual entry boundary}
\label{app:joint-root-entry}

Let $S$ be an independent uniform 32-byte contribution. Executed
preparation may choose a known 32-byte old root $r$ and the response
core using the public oracle, but has no other raw use of $S$. Consider
one original Extract/Expand,
\[
 P=F(r,S),\qquad R=F(P,I_R\Vert\mathtt{01}),
\]
with the actual 97-byte information string from
Section~\ref{sec:joint-continuation}. Repeats reuse this result.
The value $P$ has only this Expand use; $S,P$ have no other output
channel. The entire later protocol is generated from $R$.

First apply the single hidden-message column argument to salt $r$.
Its ideal value is $V(\Gamma(r))$. Even when $r$ depends on the public
oracle prefix, this first table value is fresh independent uniform:
conditional on each oracle and prepared state, its index is fixed and
has not been visited. Presampling a fresh PRK and withholding it until
this call gives exactly this joint distribution.

Embed a hidden row key at that PRK and answer the one 98-byte Expand
message. Its ideal reply is a fresh uniform root. It can again be
presampled and withheld until creation, independent of the common
oracle and all earlier state. All later confirmations, signatures,
traffic and successor operations are recomputed from this root.
The proof therefore supplies the initialization for
Theorem~\ref{thm:joint-message-leaf} to the entire interactive
continuation.

If $B$ bounds all ordinary compression work other than this derivation
and the displayed public salt-state computations, conservative budgets
are $q_{\rm ext}\le B+5$ and $q_{\rm exp}\le B+2$.
The former includes the five actual root-Expand compression calls; the
latter includes the two public states in $\Gamma(r)$.
Since $\ell(98)=3$, the resulting bit error is
\begin{equation}
\begin{split}
 \Delta_{\rm init}\le\min\{1,\;&
 \Phi_{256}(3,B+8)+E_{F_1}(1,B+7)\\
 &+\Phi_{256}(2,B+5)+E_{\rm NMAC}(1,B+2,3)\}.
\end{split}
 \label{eq:appendix-root-initialization}
\end{equation}
All terms count the complete continuation in their respective reductions.
The bound retains the published asymptotic constants and their stated
parameter regimes.

The downstream length budgets are equally explicit: confirmation
information lengths $62,58,60$ give row messages $63,59,61$; the
155-byte traffic information gives $156,188,188,188$; and the
56-byte first-step information gives $57,89$. Thus the shared root row
has $\ell(32)=2$, the traffic rows have maximum $\ell(188)=5$, and the
selected chain row has $\ell(89)=3$. These counts include the previous
HKDF output block where it is part of the actual HMAC message.
